\documentclass[11pt]{article}
\usepackage[margin=1in]{geometry}
\usepackage{amsmath,amssymb,amsthm}
\usepackage[colorlinks=true,urlcolor=blue,linkcolor=blue,citecolor=blue]{hyperref}
\newtheorem{theorem}{Theorem}
\newtheorem{lemma}[theorem]{Lemma}
\newtheorem{corollary}[theorem]{Corollary}
\newtheorem{proposition}[theorem]{Proposition}
\newcommand{\R}{\mathbb R}
\newcommand{\norm}[1]{\left\lVert#1\right\rVert}
\newcommand{\range}{\operatorname{range}}
\newcommand{\rank}{\operatorname{rank}}
\newcommand{\tO}{\widetilde O}
\newcommand{\MM}{T_{\rm MM}}
\newcommand{\Ac}{\mathcal A}
\newcommand{\Jc}{\mathcal J}
\title{Hierarchical absolute pencils for block spectral\\
approximation of factored products}
\author{Working manuscript: exact-arithmetic proof; priority under review}
\date{November 14, 2026}
\begin{document}
\maketitle
\begin{abstract}
We refine a coarse absolute-value preconditioner using finer Gram
proxies that are supplied only through sketched-factor actions.
The actual finite preconditioners satisfy a relaxed quadratic invariant
\(F-\kappa KF^{-1}K\preceq ChD\), with fixed \(\kappa>1\).
It preserves the \(\alpha h\) transfer of a strong block residual,
although the ordinary quadratic defect need not be \(O(hD)\).
Each refinement uses \(\widetilde O(H/h)\) fine-proxy actions.
An exact carried-preimage identity implements the physical output lift
without inverting an implicit polynomial preconditioner.

All depth dependence can be unfolded. For fixed
\(1/8<d<(\sqrt2-1)/2\), \(n=p=N=m^{1+d+o(1)}\),
\(\theta=(m/N)^{5/2}\), and every \(r\leq\sqrt m\), a scalar
square-root schedule of depth \(O(\log\log m)\) gives total
explicit-output time \(N^{2+o(1)}\), with gap-free error
\((1+\theta)\sigma_{r+1}(AB)\). Confidence is constant or inverse
polynomial. The bound includes all preprocessing, range algebra,
and output entries. The model counts exact real arithmetic, positive
scalar square roots, comparisons, and Gaussian samples; it makes no
stable finite-word claim. The combined theorem's global priority remains
under review.
\end{abstract}

\section{Statements and arithmetic model}
Let \(A\in\R^{n\times m}\), \(B\in\R^{m\times p}\),
\(m\leq\min(n,p)\), \(M=AB\), and \(I=m(n+p)\).
Unless specified otherwise, \(1\leq r\leq m\) is an integer.
Eigenvalues and singular values past a matrix's dimension are zero.
We use a feasible square multiplication exponent \(\omega>2\), with
fixed slack \(\zeta>0\); \(\omega\leq9/4\) is available over the
reals \cite{amanov}. Multiplication costs are chosen symmetrically in
their three dimensions, using the matrix multiplication tensor. Define
\begin{align*}
 \Ac(b)&=\MM(n,m,b)+\MM(p,m,b),\\
 s_A(h)&=\min\{n,\tO(m/h^2)\},\quad
 s_B(h)=\min\{p,\tO(m/h^2)\},\\
 \Jc_h(b)&=\MM(s_A(h),m,b)+\MM(s_B(h),m,b),\\
 G_h&=\MM(m,s_A(h),m)+\MM(m,s_B(h),m).
\end{align*}
The row counts include their allocated failure probabilities. Use the
identity embedding when a row bound reaches the ambient dimension.
Thus this fallback retains the original factors as operators and its
cost is included explicitly. Forming and storing an \(h\)-proxy costs
\(\tO(I)\); \(G_h\) is paid only when its Gram is formed.

All operations are exact field operations, comparisons, positive scalar
square roots, or Gaussian samples of unit cost. Representations may have
arbitrary size. Tildes hide logarithms in dimensions, inverse supplied
accuracy, and inverse failure probability. No input-conditioned logarithm
is hidden. A subscript \(L\) on \(\tO_L\) permits dependence on a
\emph{fixed} depth only; the growing-depth result below unfolds it.

\begin{theorem}[Fixed-depth tradeoff]\label{thm:fixed}
Fix an integer \(L\geq1\),
\[
 0<\theta\leq h_L<\cdots<h_0\leq1/64,
 \qquad h_j/h_{j-1}\leq1/1024.
\]
For every fixed factor pair and \(0<\delta<1\), an algorithm returns
factors of width at most \(r\), explicitly forms their product \(C\),
and satisfies
\(\norm{M-C}_2\leq(1+\theta)\sigma_{r+1}(M)\)
with probability at least \(1-\delta\).
Let
\[
 B_j=2^{j+1},\quad C_0=36,\quad C_j=100(C_{j-1}+1),\quad
 A_L=10^6(C_L+1),\quad \alpha=\frac{\theta}{A_Lh_L}.
\]
Take a sufficient capped per-run degree
\(q=\tO(\alpha^{-1/2})\), including the logarithm of the final
Gaussian covariance condition, and put \(w_K=(2q+1)(r+1)\).
When \(w_K\leq m\), total cost is
\begin{equation}\label{eq:fixedcost}
 \tO_L\!\left(I+G_{h_0}+m^{\omega+\zeta}
 +\Ac(r+1)\sqrt{h_L/\theta}
 +\sum_{j=1}^L\Jc_{h_j}(r+1)
                    \frac{h_{j-1}}{\sqrt{h_L\theta}}
 +np+\MM(n,r,p)\right).
\end{equation}
For arbitrary \(w_K\), adding \(m w_K^2\) gives a conservative
range-algebra bound. Rank-deficient factors, repeated values, zero
products, and zero optimal tails are covered. All outcome costs are capped.
\end{theorem}

\begin{theorem}[A growing-depth endpoint family]\label{thm:family}
Fix \(1/8<d<(\sqrt2-1)/2\) and \(c_\delta>0\).
Let \(n=p=N=m^{1+d+o(1)}\),
\[
 \rho=m/N,\qquad \theta=\rho^2\sqrt\rho,\qquad
 1\leq r\leq\lfloor\sqrt m\rfloor,
 \qquad m^{-c_\delta}\leq\delta<1.
\]
There is a depth-\(O_d(\log\log m)\) instance of the finite algorithm
whose total explicit-output cost is \(N^{2+o(1)}\), with the same
width and error guarantee as Theorem~\ref{thm:fixed}.
The schedule uses only scalar square roots, field operations, and
comparisons. More explicitly, for a fixed-shape multiplication overhead
\(\chi(m)=m^{o(1)}\), the cost is bounded by
\begin{equation}\label{eq:unfolded}
 N^2\chi(m)\exp\!\left(O_{d,c_\delta}((\log\log m)^2)\right)
       +m^{\omega+\zeta}(\log m)^{O_{d,c_\delta}(1)}.
\end{equation}
Choose fixed slack \(\zeta\) below the positive core/output margin.
\end{theorem}

\section{Precisely stated ingredients from the companion}
We use the following finite results from the independently audited
companion \cite{companion}. They are recalled to specify exactly what
is imported; all new hierarchical obligations are proved here.

Fast embeddings \cite{cnw} and the companion's product-support reduction
give explicit \(W,L_c,Q_0\) and operator identities, on the common
coarse event,
\begin{equation}\label{eq:support}
 \begin{gathered}
 M=Z\mathcal B,\qquad Z=AW,\qquad\mathcal B=L_cB,\\
 G=Z^\top Z,\qquad Q=\mathcal B\mathcal B^\top,\\
 \frac I{1+h_0}\preceq G\preceq\frac I{1-h_0},\qquad
 (1-h_0)Q\preceq Q_0\preceq(1+h_0)Q.
 \end{gathered}
\end{equation}
Its dimension \(k\) equals \(\rank M\) on the event; the computed
dimension never exceeds \(\rank M\) on any outcome. \(Z\) is injective
and \(Q_0\) is positive definite on every computed nonzero support.
The reduction costs \(\tO(I+G_{h_0}+m^{\omega+\zeta})\).
If the computed support is zero, return zero. If its dimension is at most
\(r\), project \(M\) onto the full range of \(Z\); this is exact on
the coarse event. Hence the remaining analysis has \(r<k\).

Use an even, fixed Neumann polynomial
\[
 B_m=\sum_{i=0}^{l_m}(I-G)^i,\qquad
 \varepsilon_m=\theta/10^9,
 \qquad(1-\varepsilon_m)G^{-1}\preceq B_m
                   \preceq(1+\varepsilon_m)G^{-1}.
\]
Its degree is \(O(\log(1/\theta))\); even degree makes it positive
for every positive \(G\), even off the event. Fix throughout all layers
and thresholds
\begin{equation}\label{eq:truepencil}
 K(t)=tB_m-Q,\quad D(t)=tI+Q_0,\quad
 \widehat H=B_m^{-1/2}QB_m^{-1/2}.
\end{equation}
The analytical roots are not computed. With \(H=G^{1/2}QG^{1/2}\),
\begin{equation}\label{eq:masscompare}
 \lambda_i(H)/(1+\varepsilon_m)
 \leq\lambda_i(\widehat H)
 \leq\lambda_i(H)/(1-\varepsilon_m).
\end{equation}
The eigenvalues of \(H\) are the nonzero squared singular values of \(M\).

The companion's \emph{Finite weighted defect} lemma constructs an
explicit \(F_0=J_0J_0^\top\) and \(J_0^{-1}\), at every \(t>0\),
in \(\tO(m^{\omega+\zeta})\), with
\begin{equation}\label{eq:base}
 8h_0D\preceq F_0\preceq2D,\quad -F_0\prec K\prec F_0,
 \quad0\preceq F_0-KF_0^{-1}K\preceq36h_0D.
\end{equation}
Set \(R_0=J_0^{-\top}\).
It computes \(S_0=(tI-Q_0)(tI+Q_0)^{-1}\), a finite upper Newton
approximation to \(\sqrt{S_0^2+64h_0^2I}\), and its regularized
coarse metric. This is finite field/Cholesky algebra, not a spectral
oracle. The coarse core remains positive on every outcome.

The companion's \emph{Coefficient ratio and shifted inertia} lemma,
using fast characteristic polynomials \cite{np}, proves
\(v_j/j\leq e_j/e_{j-1}\leq(k-j+1)v_j\) for PSD spectra, and
computes exact shifted inertia in \(\tO(m^{\omega+\zeta})\), including
zero and repeated roots. Thus for \(j=r+1\) and
\(\varrho=e_j(Q_0)/e_{j-1}(Q_0)\),
\begin{equation}\label{eq:init}
 [\ell,u]=[\varrho/(4k),4k\varrho]
 \quad\hbox{contains}\quad\lambda_*:=\lambda_j(\widehat H)>0.
\end{equation}
Its ratio is \(16k^2\); no leading-to-target scale search is needed.

We also import the companion's \emph{Safe threshold oracle}. Its proof
applies to \emph{any} symmetric \(C\) with \(-I\prec C\prec I\),
\(T=I-C\), and an inertia-equivalent pencil. Musco--Musco's per-vector
Theorem~12 \cite{mm}, applied to target \(j\), gives a compressed
\(T^2\) Ritz value \(\beta_j\leq\lambda_j(T)^2\) on every run and
\(\beta_j\geq(1-\alpha/100)\lambda_j(T)^2\) with probability
\(99/100\). Independent repetitions and exact shifted-inertia tests at
\(1\) and \(1-2\alpha/100\) return one of:
\[
 \textit{below}:\ \mu_j(C)<0,\qquad
 \textit{above}:\ \mu_j(C)\geq0,\qquad
 \textit{near}:\ |\mu_j(C)|\leq\alpha.
\]
Here \(\mu_i\) is ordered increasingly. A below decision is always safe
on the common event; failure at most \(\beta\) for the whole query
requires only logarithmic repetitions. The degree is
\(\tO(\alpha^{-1/2})\). No numerical Ritz extraction, root isolation,
or coefficient powering is used.

Finally, deterministic CUP rank revelation \cite{jps}, followed by
exact independent-column Gram--Cholesky algebra, supplies range bases
in \(\tO(m^{\omega+\zeta})\) whenever both dimensions are at most
\(m\). The physical thin-basis cost is bounded by
\(\Ac(r)+m^{\omega+\zeta}\) using tensor symmetry and padding.
These algorithms permit arbitrary column dependencies.

\section{Fine inverse masses use only fine factors}
Store finer factors \(A_j=S_jA\), \(B_j^{\rm fac}=BT_j^\top\)
without forming their Grams. On their simultaneous Loewner events,
\[
 G_j=W^\top A_j^\top A_jW\in[(1-h_j)G,(1+h_j)G],\qquad
 Q_j=L_cB_j^{\rm fac}(B_j^{\rm fac})^\top L_c^\top
                         \in[(1-h_j)Q,(1+h_j)Q].
\]
Coarse whitening gives \(\norm{I-G_j}_2\leq2h_0<1\).
An even Neumann polynomial \(\widetilde B_j\) approximates
\(G_j^{-1}\) relatively to \(h_j/100\) in
\(O(\log(1/h_j))\) fine-factor actions. Define
\(K_j(t)=t\widetilde B_j-Q_j\). Inverse Loewner comparison and
\eqref{eq:support} imply, uniformly for every \(t>0\),
\begin{equation}\label{eq:proxy}
 -8h_jD\preceq E_j:=K_j-K\preceq8h_jD.
\end{equation}
Indeed the two inverse masses differ by at most a small constant times
\(h_jI\); the right proxy error is at most
\(h_jQ\preceq h_jQ_0/(1-h_0)\). The constant eight is conservative.
All inverse masses are fixed polynomials throughout the threshold search.
Their actions cost \(\tO(\Jc_{h_j}(b))\) on a width-\(b\) block.
They use the stored fine factors, including an identity fallback when
chosen, with no finer support calculation.

\section{A finite relaxed invariant}
The stronger ordinary defect bound in \eqref{eq:base} will be replaced
by
\begin{equation}\label{eq:invariant}
 8h_jD\preceq F_j\preceq B_jD,\qquad -F_j\prec K\prec F_j,
 \qquad F_j-\kappa_jKF_j^{-1}K\preceq C_jh_jD,
\end{equation}
where \(\kappa_0=1\) and \(\kappa_j=5\) for \(j\geq1\).
The last left side need not be positive. This distinction permits a
hierarchy without asserting a false \(O(h_jD)\) ordinary defect.

\begin{lemma}[Relative perturbation]\label{lem:relative}
Let \(F_*,\bar F\) be positive definite, \(0\leq\delta<1\), and \(\kappa\geq0\).
If \((1-\delta)F_*\preceq\bar F\preceq(1+\delta)F_*\) and
\(F_*-\kappa KF_*^{-1}K\preceq ChD\), then
\[
 \bar F-\kappa(1+\delta)^2K\bar F^{-1}K
 \preceq(1+\delta)(F_*-\kappa KF_*^{-1}K)
 \preceq(1+\delta)ChD.
\]
No matrices in this statement need commute.
\end{lemma}
\begin{proof}
Use \(\bar F\preceq(1+\delta)F_*\) and
\(\bar F^{-1}\succeq(1+\delta)^{-1}F_*^{-1}\), then apply
congruence by \(K\).
\end{proof}

\begin{lemma}[Finite inverse-root polynomial]\label{lem:poly}
For \(0<\eta\leq1/8\), \(0<\tau<1/4\), and
\(f_\eta(x)=\sqrt{x^2+\eta^2}+\eta\), a finite polynomial \(p\)
of degree \(O(\eta^{-1}\log(1/(\tau\eta)))\) satisfies
\[
 |p(x)\sqrt{f_\eta(x)}-1|\leq\tau\quad(-1\leq x\leq1).
\]
Its coefficients use field operations and positive scalar square roots
only, in \(O(n\log n)\) work for a dyadic degree bound \(n\).
\end{lemma}
\begin{proof}
Let \(g(z)=(\sqrt{z^2+\eta^2}+\eta)^{-1/2}\), with positive real
branches. On the Bernstein ellipse of parameter
\(\varrho=1+\eta/4\), the imaginary part has magnitude at most
\(\eta/4\). Hence \(\operatorname{Re}(z^2+\eta^2)\geq15\eta^2/16\);
the branches are analytic there and \(|g(z)|\leq\eta^{-1/2}\).
The Laurent representation under \(z=(w+w^{-1})/2\) and Cauchy's
bound give Chebyshev coefficient magnitudes at most
\(2\eta^{-1/2}\varrho^{-i}\). Interpolation at \(n\) Gauss--Chebyshev
nodes aliases each omitted Chebyshev polynomial to a polynomial of norm
at most one on the interval. The interpolation error is therefore at
most \(24\eta^{-3/2}\varrho^{-n}\).

Choose dyadic \(n\) by repeatedly squaring \(\varrho^{-n}\) until it
is at most \(\tau\eta\sqrt\eta/1000\). This gives the stated degree
and absolute error at most \(\tau/2\). Since
\(g(x)\geq1/\sqrt{1+2\eta}>0.8\), it gives the relative claim and
positivity of \(p\). The dyadic nodes are real parts of dyadic roots
of unity, constructed by positive half-angle square roots. Each node
value of \(g\) uses two positive scalar roots and division. The dyadic
DCT is evaluated by the usual divide-and-conquer DFT recurrence using
these roots, in \(O(n\log n)\) real-pair operations. Thus no cosine,
integration, coefficient, or matrix-root oracle is used. The Chebyshev
recurrence applies \(p\) using \(n\) matrix actions.
\end{proof}

\begin{lemma}[One finite refinement]\label{lem:refine}
Assume \eqref{eq:invariant} at accuracy \(H=h_{j-1}\), with
\(\kappa\leq5\), constants \(B,C\), and \(C\geq1\).
Let \(h=h_j\), \(\epsilon=h/H\leq1/1024\), and suppose
\eqref{eq:proxy} holds. If an actual inverse factor \(R_{j-1}\) has
\(F_{j-1}=R_{j-1}^{-\top}R_{j-1}^{-1}\), a finite polynomial
construction produces \(R_j\) satisfying \eqref{eq:invariant} with
\[
 B_j=2B,\qquad C_j=100(C+1),\qquad\kappa_j=5.
\]
It uses degree \(\tO(H/h)\) and only forward previous-level and
fine-proxy actions. In addition \(F_j\pm K\succeq55hD\).
\end{lemma}
\begin{proof}
Set \(P=(1+\epsilon)F_{j-1}\),
\(R_P=R_{j-1}/\sqrt{1+\epsilon}\), and
\(A_j=R_P^\top K_jR_P\). The error bound and
\(F_{j-1}\succeq8HD\) give \(-P\prec K_j\prec P\), hence
\(\operatorname{spec}(A_j)\subset[-1,1]\).

Write \(E=K_j-K\). Weighted normalization gives
\(EF_{j-1}^{-1}E\preceq8h^2D/H\). Quadratic Young's inequality,
with parameter \(\epsilon\), gives
\[
 KF_{j-1}^{-1}K\preceq(1+\epsilon)K_jF_{j-1}^{-1}K_j
 +(1+1/\epsilon)EF_{j-1}^{-1}E.
\]
Consequently
\begin{equation}\label{eq:young}
 P-\kappa_PK_jP^{-1}K_j\preceq C_PH D,\quad
 \kappa_P=\kappa(1+\epsilon)^3\leq6,\quad C_P\leq2(C+1).
\end{equation}
This controls noncommuting cross terms; no commutation with \(D\) or
\(E\) is assumed.

Put \(\eta=8\epsilon\) and define, for analysis only,
\(F_*=R_P^{-\top}f_\eta(A_j)R_P^{-1}\).
Since \(f_\eta\geq|x|+\eta\), its minimum is \(2\eta\), and its
maximum is at most \(1+2\eta\leq5/4\),
\begin{equation}\label{eq:idealbounds}
 F_*\succeq128hD,\quad F_*\preceq(1+2\eta)(1+\epsilon)BD,
 \quad F_*\pm K\succeq56hD.
\end{equation}
The scalar identity \(f_\eta-x^2/f_\eta=2\eta\) gives exactly
\(F_*-K_jF_*^{-1}K_j=2\eta P\).
By \eqref{eq:young} and \(F_*\preceq(5/4)P\),
\[
 F_*-K_jF_*^{-1}K_j
 \preceq2\eta\kappa_P(5/4)K_jF_*^{-1}K_j+16C_PhD.
\]
The first coefficient is below one because
\(\epsilon\leq1/1024\). Hence
\(F_*-2K_jF_*^{-1}K_j\preceq16C_PhD\).
Young's inequality with parameter one and
\(EF_*^{-1}E\preceq hD/2\) now imply
\begin{equation}\label{eq:idealrelax}
 F_*-4KF_*^{-1}K\preceq(16C_P+2)hD.
\end{equation}

Apply Lemma~\ref{lem:poly} with
\(\tau=h/[10^6(1+B)]\), and set
\[
 R_j=R_Pp_j(A_j),\qquad F_j=R_j^{-\top}R_j^{-1}.
\]
Only \(R_j\) and its transpose are applied. The analytical metric
satisfies \((1-4\tau)F_*\preceq F_j\preceq(1+4\tau)F_*\).
Lemma~\ref{lem:relative} turns \eqref{eq:idealrelax} into the claimed
\(\kappa_j=5\), \(C_j=100(C+1)\) inequality. The relative error,
\eqref{eq:idealbounds}, and the chosen \(\tau\) preserve the floor,
the \(55hD\) margins, and the upper bound \(2BD\).
All constants are uniform in the level; \(C\) enters linearly, and
\(B\) enters only the scalar tolerance. No fine metric is formed.
\end{proof}

Starting at \eqref{eq:base}, Lemma~\ref{lem:refine} proves the actual
finite invariant at every level with the explicit theorem constants.
The construction is uniform over all positive thresholds, including
singular pencils. It uses no fine spectral-norm refinement or fine Gram.

\section{Threshold search and physical transfer}
At a fixed threshold set
\(F=F_L\), \(R=R_L\), \(C_F=R^\top KR\), \(T=I-C_F\).
Domination gives \(0\prec T\prec2I\). The safe threshold oracle from
the companion applies literally to this symmetric operator. The congruent
matrix \(C_F\) has the inertia of \(K=tB_m-Q\), hence counts
eigenvalues of \(\widehat H\) above the threshold.
Update \(\ell=t\) on below and \(u=t\) on above. At a near result,
\(|\mu_j(C_F)|\leq\alpha\).

\begin{lemma}[Relaxed min--max and strong transfer]\label{lem:transfer}
Suppose \(-F\prec K\prec F\) and
\(F-\kappa KF^{-1}K\preceq ChD\), with \(1\leq\kappa\leq5\).
Assume \(0<\alpha\leq1/(100\kappa)\) and \(7C\alpha h<1/4\),
as holds for the theorem's parameters at \(C=C_L,h=h_L\).
For the theorem's \(\alpha\), a near result implies
\[
 \lambda_*\geq\frac{1-x}{1+x}t,\qquad
 x\leq4C\alpha h.
\]
If \(U\) is an orthonormal basis satisfying
\(\norm{(I-UU^\top)T}_2\leq1+3\alpha\), then
\(\mathcal S=\range(ZR^{-\top}U)\) satisfies
\begin{equation}\label{eq:physical}
 \norm{(I-\Pi_{\mathcal S})M}_2^2
 \leq\frac{1+u+\varepsilon_m}{1-u}\,t,
 \qquad u=7C\alpha h.
\end{equation}
\end{lemma}
\begin{proof}
Put \(\beta_\alpha=\alpha/(1-\kappa\alpha^2)\leq2\alpha\).
For \(-1\leq c\leq\alpha\),
\(c\leq\beta_\alpha(1-\kappa c^2)\). To check this, the convex
quadratic \(\beta_\alpha\kappa c^2+c-\beta_\alpha\) vanishes at
\(c=\alpha\) and is nonpositive at \(c=-1\).
Apply it on the bottom-\(j\) eigenspace, undo the congruence, and use
the relaxed invariant. On a dimension-\(j\) subspace,
\(K\preceq C\beta_\alpha hD\preceq2C\beta_\alpha h(tB_m+Q)\).
Generalized min--max proves the first assertion.

Let \(b=(1+3\alpha)^2-1\leq7\alpha\). Symmetry and the strong
action norm imply \(w^\top(C_F^2-2C_F)w\leq b\norm w^2\) for
\(w\perp U\). Substitute \(w=R^{-1}y\). Then
\[
 2y^\top Ky\geq y^\top KF^{-1}Ky-b y^\top Fy
 \geq(1-\kappa b)y^\top KF^{-1}Ky-bCh\,y^\top Dy
 \geq-bCh\,y^\top Dy.
\]
For physical \(x\perp\mathcal S\), \(y=Z^\top x\) has this
orthogonality and satisfies
\(y^\top G^{-1}y\leq\norm x^2\),
\(\norm y^2\leq2\norm x^2\), and
\(y^\top Q_0y\leq2x^\top MM^\top x\).
Using \(B_m\preceq(1+\varepsilon_m)G^{-1}\) gives
\[
 t\norm x^2-x^\top MM^\top x
 \geq-(Cbh+\varepsilon_m)t\norm x^2
       -Cbh\,x^\top MM^\top x.
\]
Taking the largest quotient proves \eqref{eq:physical}.
\end{proof}

If no near result occurs, stop when \(u-\ell\leq\alpha h_L\ell\)
and take the final \(t=u\). Inertia gives \(\lambda_j(T)\leq1\)
and \(t\leq(1+\alpha h_L)\lambda_*\). A near result gives
\(\lambda_j(T)\leq1+\alpha\). In both cases the final strong filter
is applicable. All thresholds are positive; the coarse core has no poles.
An outcome-independent bound on threshold calls is
\begin{equation}\label{eq:Qcap}
 Q_{\max}=\left\lceil\log_2(16m^2 A_L/\theta)\right\rceil+1.
\end{equation}
Scalar halving finds a sufficient integer without a logarithm oracle.

With \(C=C_L\), \(h=h_L\), the choices in the theorem give
\(u\leq7\theta/10^6\) and \(x\leq4\theta/10^6\).
The physical factor is at most \(1+3u+2\varepsilon_m\), the near
threshold costs at most \(1+4x\), and mass comparison costs at most
\(1+2\varepsilon_m\). Their product is below \((1+\theta)^2\).
This calculation keeps \(\varepsilon_m=\theta/10^9\) separate from
\(\alpha h_L=\theta/A_L\), including at growing depth.

\section{Exact carried preimages implement the dual output}
The physical lift in Lemma~\ref{lem:transfer} contains \(R_L^{-\top}\).
It cannot be replaced by an uncharged approximate inverse: that would
change the exact dual complement. We now implement the lift using only
forward fine polynomials.

Let
\[
 S=R_0^{-1}R_L,\qquad A_0=R_0^\top KR_0,\qquad
 T=I-S^\top A_0S,\qquad V_0=I-A_0SS^\top.
\]
The explicit coarse matrices apply \(R_0^{-1}\) and its transpose.
Although the fine polynomial factors need not commute,
\begin{equation}\label{eq:Scond}
\kappa(S)\leq U_S:=2^L\sqrt{h_0/h_L}
\end{equation}
Indeed each positive polynomial factor has condition at most
\(2\sqrt{h_{j-1}/h_j}\), by its relative approximation and the
minimum and maximum of \(f_\eta\); multiply the conditions.
Scalar inflations have condition one.

Use a fresh standard Gaussian \(\Xi\in\R^{k\times r}\) and the
start \(\Omega=S^\top\Xi\) for the final filter. The companion's
filter polynomial is
\[
 \phi_v(x)=\frac{{\rm Cheb}_v(2x/(1+\alpha)-1)}
 {{\rm Cheb}_v(2(1+2\alpha)/(1+\alpha)-1)}.
\]
Set the computable bound \(R_G=(100kU_S/\beta)^3\). Increase the
denominator by the scalar recurrence until it is at least \(2R_G/\alpha\).
The intertwining identity
\(TS^\top=S^\top V_0\) proves exactly
\begin{equation}\label{eq:preimage}
 Y=\phi_v(T)\Omega=S^\top W_o,\qquad
 W_o=\phi_v(V_0)\Xi,\qquad
 R_L^{-\top}Y=R_0^{-\top}W_o.
\end{equation}
Thus compute \(W_o\) by forward polynomial recurrence and output the
physical span \(\range(ZJ_0W_o)\). No inverse fine polynomial is applied.
The carried operator \(V_0\) may be nonnormal; the proof is the exact
polynomial identity, not a nonsymmetric Lanczos assertion.

For completeness, the Gaussian filter capture remains valid with this
start. In a fixed \(T\) eigenbasis, let the head be the eigenvalues
strictly above \(1+\alpha\), of dimension at most \(r\). Its Gaussian
row covariance has eigenvalues between \(\sigma_{\min}(S)^2\) and
\(\sigma_{\max}(S)^2\). A square head subblock has least singular
value at least \(\sigma_{\min}(S)\) times that of a standard Gaussian
square matrix. The elementary column-distance/union bound and a
Frobenius Markov bound for the tail give, with failure at most \(\beta\),
\[
 \norm{\Omega_{\rm tail}\Omega_{\rm head}^\dagger}_2
\leq R_G.
\]
No head/tail independence is required. Strong polynomial values at
eigenvalues at least \(1+2\alpha\) are at least one; weak head
coordinates are forced to zero by the analytical head right inverse.
Choosing the denominator at least \(2R_G/\alpha\)
gives strong-space error at most \(\alpha/2\). The remaining spectrum
is below \(1+2\alpha\), and \(\norm T\leq2\), so
\(\norm{(I-\Pi_{\range Y})T}\leq1+3\alpha\).
The required degree is
\(O(\alpha^{-1/2}\log(kU_S/(\alpha\beta)))\).
This proves the strong certificate for the actual physical space in
\eqref{eq:preimage}. Its basis may have width less than \(r\).

Select independent columns of \(W_o\), lift them by \(ZJ_0\), and
orthonormalize by exact Gram--Cholesky. Compute \(M^\top U_p\) through
the original factors and expand \(U_p(M^\top U_p)^\top\).
These are exactly the charged physical range and output operations.

\section{Fixed-depth action cost and probability}
Write \(c_j(b)\) for an action of \(R_j\) or \(R_j^\top\) on a
width-\(b\) block at one fixed threshold. Lemma~\ref{lem:poly} and
the recursive definition give
\[
 c_j(b)=\tO_L\!\left(\frac{h_{j-1}}{h_j}
                   [c_{j-1}(b)+\Jc_{h_j}(b)]\right).
\]
Only the explicit coarse metric is factored at each threshold.
Unrolling the ratios telescopes to \(h_0/h_L\). The true \(K\)
appears once in each outer \(T\) action; the inverse masses inside
the finer metrics use fine factors only. Consequently the outer action
cost, at \(q=\tO_L(\sqrt{h_L/\theta})\), is exactly the two action
terms in \eqref{eq:fixedcost}. The same is true of \(V_0\) for the
final carried filter. No intermediate norm searches are needed.

The raw outer Krylov basis has at most \(w_K\) columns, including
dependent ones. Under \(w_K\leq m\), CUP and independent-column
Gram--Cholesky cost \(\tO(m^{\omega+\zeta})\), as do its compressed
Gram and exact sign counts. If wider, conventional rank revelation and
orthogonalization cost at most \(m w_K^2\). Applying \(T\) to its
range basis in width-\((r+1)\) chunks costs the same order as generating
the raw basis. Gaussian repetitions have separate bases. Internal fine
polynomial recurrences retain only constant-width blocks.

Allocate \(\delta/2\) over the \(2(L+1)\) embedding events, and
\(\beta=\delta/[4(Q_{\max}+1)]\) to each conditional threshold query
and the final fresh filter. All preconditioner polynomials are deterministic
given the embeddings; there are no hidden randomized recursive calls.
Conditional union bounds give failure at most \(\delta\).
All layers, degrees, repetitions, rank checks, and dimensions are capped
from parameters and computed support. Failed fine embeddings can make
an effective factor singular, but the executed algorithm takes only
forward polynomials and \eqref{eq:preimage}; no inverse fine factor is
taken on these outcomes. Coarse Cholesky and coefficient denominators
are checked, and an invalid check returns zero factors within the cap.
The rank-zero and zero-tail branches were handled before threshold search.
This proves Theorem~\ref{thm:fixed}.

\section{Growing depth with every dependence unfolded}
We now prove Theorem~\ref{thm:family}; no \(\tO_L\) assertion is used.
Set \(h_0=1/1024\), \(h_*=h_0\sqrt\rho\), and start at \(h_0\).
Accept \(h_{j+1}=\sqrt{h_jh_*}\) only while
\(h_{j+1}\leq h_j/1024\). This is a finite scalar-computable schedule.

If at least one level is accepted, write \(g_j=\log(h_j/h_*)\) for
analysis. Acceptance means \(g_j\geq2\log1024\); each accepted step
halves \(g_j\). Thus the terminal value lies in
\([\log1024,2\log1024)\), and
\begin{equation}\label{eq:schedule}
 \sqrt\rho\leq h_L<1024\sqrt\rho,\quad
 L=O_d(\log\log m),\quad2^L=\Theta_d(\log m),\quad
 h_j^2=h_{j-1}h_0\sqrt\rho.
\end{equation}
The logarithms describe the finite loop; none is needed to execute it.
If no level is accepted, then \(\rho>1024^{-4}\). For fixed positive
\(d\) this gives only finitely many \(m\). Use the companion's verified
one-level algorithm at internal accuracy \(\min(\theta,h_0)\), with
its conservative cost. It preserves width at most \(r\); merely
returning the exact product could violate that requirement. The finite
exception is absorbed in the family constant.

The explicit invariant constants satisfy
\[
 B_L=2^{L+1}=\operatorname{polylog}(m),\qquad
 C_L=(36+100/99)100^L-100/99=\operatorname{polylog}(m).
\]
The refinement ratio restriction is uniform in \(C_j\): \(\kappa_j\)
resets to five at every layer. There is no hidden depth-dependent
smallness requirement. Since \(\tau_j=h_j/[10^6(1+B_{j-1})]\),
the interpolation degrees obey the common bound
\begin{equation}\label{eq:degreeuniform}
 n_j\leq C_d(h_{j-1}/h_j)\log(2m).
\end{equation}
The fine inverse-mass degree is also \(O_d(\log m)\).

Put \(\ell_m=\log(2m)\). For leaf product costs \(J_j^{\rm leaf}(b)\),
an explicit uniform recurrence is
\begin{equation}\label{eq:recuniform}
 c_j(b)\leq C\ell_m^2\frac{h_{j-1}}{h_j}
                 [c_{j-1}(b)+J_j^{\rm leaf}(b)].
\end{equation}
One factor \(\ell_m\) covers interpolation and the other fine inverse
mass; this conservative bound also covers scalar recurrence work.
Unrolling multiplies leaf sums by at most
\begin{equation}\label{eq:E}
 E(m)=(C\ell_m^2)^L
        =\exp(O_d((\log\log m)^2))=m^{o(1)}.
\end{equation}
\(B_j\) enters degrees only logarithmically; \(C_L\) enters outer
accuracy through the polynomial factor \(A_L\). They do not exponentiate
\eqref{eq:E} a second time. The largest ratio \(h_{j-1}/h_j\) is
\(\rho^{-1/4}\), so all scalar dyadic interpolation preparations cost
at most \(O_d(L\rho^{-1/4}(\log m)^2)\). They may be reused at every
threshold and are below the final bound.

\paragraph{A uniform fixed kernel.}
Use the actual fixed-shape bound
\[
 \MM(m,m,\lceil\sqrt m\rceil)\leq m^2\chi(m),
 \qquad\chi(m)=m^{o(1)},
\]
from \cite{amanov}. Split each varying row dimension into \(m\)-row
blocks and pad rank width to at most two width-\(\lceil\sqrt m\rceil\)
blocks. Transpose actions follow by tensor symmetry and summing inner
blocks. A factor with \(s\geq m\) rows costs at most
\(Cs m\chi(m)\) per block action, uniformly across every layer.
Therefore
\begin{equation}\label{eq:kernels}
 \Ac(r+1)\leq CmN\chi(m),\quad
 J_j^{\rm leaf}(r+1)
   \leq Cm^2h_j^{-2}\chi(m)(\log m)^\gamma
\end{equation}
for a fixed embedding-log exponent \(\gamma\).
The identity fallback only decreases the stored row bound.
There is no varying-shape, nonuniform rectangular-exponent assumption.
The \(\chi(m)\) factor occurs once per expanded leaf, rather than
being raised to \(L\), because \eqref{eq:recuniform} is linear in costs.

All stored sketches have at most \(N\) rows. Their combined storage
is \(O(LmN)\), and application costs \(mN\) times fixed logarithmic
factors. No fine Gram is formed, even if the final levels use identity
factors. Coarse sketch rows are \(m\operatorname{polylog}m\); their
Grams and support cost \(m^{\omega+\zeta}\operatorname{polylog}m\).

The exact cancellations in \eqref{eq:schedule}, with
\(\theta=\rho^{5/2}\), give
\begin{equation}\label{eq:cancel}
 mN\sqrt{h_L/\theta}\leq32N^2,\qquad
 \frac{m^2h_{j-1}}{h_j^2\sqrt{h_L\theta}}
 =\frac{m^2}{h_0\sqrt\rho\sqrt{h_L\theta}}
 \leq N^2/h_0.
\end{equation}
Every unrolled fine term is thus quadratic before the common overhead.
Combine \eqref{eq:E}--\eqref{eq:cancel}, \(\sqrt{A_L}\), and the
threshold/source logarithms to get \eqref{eq:unfolded}.
Since \(d>1/8\), choose fixed \(\zeta<2+2d-9/4\); the core term
is smaller than \(N^2\). Thus the total is \(N^{2+o(1)}\).

For clarity, the actual outer degree and raw width obey
\[
 q\leq(N/m)\operatorname{polylog}m=m^{d+o(1)},\qquad
 w_K\leq m^{1/2+d+o(1)}<m.
\]
Here \(\log\kappa(S)\leq O(L)+\tfrac12\log(h_0/h_L)=O_d(\log m)\),
so elliptical capture adds only a logarithm. Every core range/sign-count
operation uses the fast bound. Gaussian repetitions never concatenate
bases; inner polynomial recurrences do not enlarge the outer raw width.
\eqref{eq:Qcap} is \(O_d(\log m)\), and with inverse-polynomial
confidence all embedding/query failure logarithms are \(O(\log m)\).
The deterministic finite invariants introduce no extra failure events.

Finally, tile the \(N\)-by-\(N\) output into \(m\)-by-\(m\) blocks.
Its inner rank is at most \(\sqrt m\), so the same tensor-permuted
fixed kernel expands it in \(O(N^2\chi(m))\), including every entry.
Physical Gram--Cholesky uses the same row tiling plus \(m^{\omega+\zeta}\).
All outcome caps and the exact zero-tail branch remain valid at this
growing depth. This completes Theorem~\ref{thm:family}.

\section{An optimized total-output comparison}
The comparison uses fast factor-block actions throughout, including when
the target rank grows polynomially. No factor \(r\) is inserted into an
input-optimal block product.

\begin{corollary}\label{cor:optimized}
Put
\[
 d_* =\frac{3\sqrt5-5}{10},\qquad
 \alpha_* =\frac{3\sqrt5-5}{2},\qquad
 s_* =\frac{7-3\sqrt5}{8}=0.0364745084\ldots .
\]
For \(n=p=N\), \(m=N^{\alpha_*+o(1)}\),
\(\theta=(m/N)^{5/2}=N^{-10s_*+o(1)}\), and every
\(1\leq r\leq\lfloor\sqrt m\rfloor\), Theorem~\ref{thm:family}
gives total time \(N^{2+o(1)}\).
Both the checked ordinary block-Krylov bound and the exact-Gram bound
have exponent \(2+s_*\) at these parameters. In particular the rank
may grow through \(N^{\alpha_*/2+o(1)}\).
The comparison is of available arithmetic upper bounds, not a lower
bound for every algorithm.
\end{corollary}
\begin{proof}
For a general fixed \(d\) in Theorem~\ref{thm:family}, a width-at-most-
\(\sqrt m\) factor action costs \(m^{2+d+o(1)}\). Ordinary block
Krylov \cite{mm} uses \(\theta^{-1/2}=m^{5d/4+o(1)}\) iterations,
giving the checked exponent \(2+9d/4\) in powers of \(m\).
This includes postprocessing: its raw basis width has exponent
\(\ell=1/2+5d/4<1\). Fast tall rank revelation and range algebra cost
at most exponent
\(1+d+(5/4)\ell=13/8+41d/16<2+9d/4\).
The compressed Gram exponent
\(1+d+\ell+\ell^2/[4(1+d)]\) is smaller still.
Even the conventional cubic dense Ritz charge has exponent
\(3\ell=3/2+15d/4<2+9d/4\), since \(d<1/3\).
Full-basis factor actions can be split into width-\(r\) blocks and
cost at most the same iteration term. Final projection and expansion
also fit that term. A uniform exact-field realization of this same
scaling, without a Ritz oracle, follows from the companion's theorem
at constant distortion \(h=1/1024\); its raw-width condition also
holds because \(1/2+5d/4<1\).

Exact Gram formation has exponent
\(2+d+1/[4(1+d)]\) by \cite{amanov}; the remaining core, lifting,
and output work are smaller on the stated interval. Appendix~\ref{app:gram}
gives a realization in the same exact-field model without an SVD oracle.
Subtracting the new total exponent two after conversion to powers of
\(N\), the two comparison gaps are
\[
 g_K(d)=\frac{d}{4(1+d)},\qquad
 g_G(d)=\frac{1}{4(1+d)^2}-\frac{d}{1+d}.
\]
The first increases and the second decreases. Their intersection is
\(5d(1+d)=1\), namely \(d=d_*\), where both gaps equal \(s_*\).
The identities \(\alpha_*=1/(1+d_*)\) and
\(10s_*=5d_*/[2(1+d_*)]\) give the displayed parameters.
This endpoint saving is attained for exponent-two total time; it is
not a literal \(O(N^2)\) or strict subquadratic-construction statement.
The earlier scalar stochastic route already reaches this accuracy
envelope at fixed rank. The hierarchy extends it through growing ranks;
no dominance over that route in every regime is claimed.
\end{proof}

\begin{corollary}[Absolute error without a supplied rank]\label{cor:absolute}
In the dimensional family of Theorem~\ref{thm:family}, let \(e>0\)
be supplied and define
\[
 r_e=\min\{s\in\{0,\ldots,m\}:\sigma_{s+1}(M)\leq e/(1+\theta)\}.
\]
Without being supplied \(r_e\), an algorithm returns explicit error at
most \(e\) and width at most \(2r_e\), with probability at least
\(1-\delta\). On that same event its total time is \(N^{2+o(1)}\)
whenever \(r_e\leq\lfloor\sqrt m\rfloor/2\). Other outcomes have
a finite dimension-dependent cap; this adaptive cost is a high-probability
statement.
\end{corollary}
\begin{proof}
Fix \(t=e^2/(1+\theta/4)\) and use the companion's rank sequence
\(0,1,2,4,\ldots\). Construct the present hierarchical metric at that
one threshold.
At each attempt, take the full-support projection branch first if
\(r\) reaches the computed support dimension, before requesting \(j=r+1\).
For smaller widths, advance if any compressed squared quantile exceeds
one. If all are at most one, any successful per-vector repetition gives
\(\lambda_{r+1}(T)\leq1+\alpha\). The fresh carried filter and
\eqref{eq:physical} give error at most \(e\), since
\(3u+2\varepsilon_m<\theta/4\). At rank zero the accepted norm
bound itself gives the strong certificate for the empty output space.
The represented tail at \(r_e+1\) is strictly below \(t\), by
\eqref{eq:masscompare}. Thus no attempt at or above \(r_e\) can
advance on the common embedding event, even if its per-vector runs
fail. The stopping rank is at most \(2r_e\).

If the next attempt would exceed \(\lfloor\sqrt m\rfloor\), use
the companion's globally capped absolute-error algorithm, with an
independent half-budget for failure. It also gives width at most
\(2r_e\); an internal accuracy smaller than \(\theta\), if needed,
only decreases its slack rank. On the stated small-\(r_e\) event this
fallback is never reached. On this small-\(r_e\) event, every hierarchical
block is at most \(\sqrt m+1\),
and the logarithmic number of attempts is absorbed in \eqref{eq:unfolded}.
Conditional failure budgets, including a reserved fallback budget, give
one common event for error, width, and time. The no-level finite-size
branch uses that same companion algorithm.
\end{proof}

\section{A two-level growing-rank example and scope}
\begin{corollary}\label{cor:twolevel}
Let \(n=p=N=m^{59/50+o(1)}\),
\[
 h_0=m^{-7/100},\quad h_1=m^{-1/12},\quad
 \theta=m^{-11/25},\quad1\leq r\leq\lceil m^{14/25}\rceil.
\]
For constant confidence, the two-level algorithm constructs the
approximation in \(N^{6724/3363+o(1)}\) and explicitly outputs it in
total \(N^{2+o(1)}\).
\end{corollary}
\begin{proof}
The ratio condition holds eventually, and the raw width exponent is
\(443/600<1\). At the endpoint rank, the full rectangular formula in
\cite{amanov} gives the following exponents in powers of \(m\):
\[
\begin{array}{c|c}
 G_{h_0}&3362/1425\\
 \Ac(r+1)\sqrt{h_1/\theta}&29227553/12390000\\
 \Jc_{h_1}(r+1)h_0/\sqrt{h_1\theta}&2831/1200\\
 m^\omega&9/4
\end{array}
\]
The fine block cost itself has exponent \(867/400\); it is slightly
outside the input-optimal range and is not charged just as its input
size. Coarse Gram is largest, below output \(59/25\), and exceeds
the fine term by \(1/7600\). Fine sketch rows have exponent
\(7/6<59/50\), so neither sketch is an identity fallback here.
All strict margins accommodate logarithms and small multiplication slack.
Since \(r=N^{28/59+o(1)}<N^{1/2}\), expansion has exponent two.
\end{proof}

Power-law parameters in this example may be rounded within fixed factors
to dyadic values using integer-power comparisons. No general root oracle
is needed to select them. This example expands the available one-level
accuracy tradeoff at growing rank: even rank-one true-factor work at
\(\theta=m^{-11/25}\) would require a one-level distortion exponent
at least \(2/25\) to fit \(N^2\), whose sketched-Gram exponent is
\(54/25+25/116>59/25\). This compares checked upper-bound tradeoffs;
it is not a lower bound for every one-level algorithm.

The \(H/h\) refinement degree is meaningful. A uniform polynomial
inverse congruence satisfying \(|x|p(x)^2\leq1\) on \([-1,1]\),
\(p(0)^2\geq1/(C\epsilon)\), and \(C\epsilon\leq1/16\) has degree
at least \(1/(20C\epsilon)\), hence \(\Omega(1/\epsilon)\) for fixed \(C\).
Indeed if \(M_p=\max|p|\), its maximizer lies within \(1/M_p^2\)
of zero; at \(4/M_p^2\) the value has magnitude at most \(M_p/2\).
Bernstein's interior derivative bound gives
\(\deg p\geq M_p^2/20\). At a zero threshold the relaxed invariant
forces precisely this large central value, with \(\epsilon=h/H\).
Thus the old positive inverse-root hierarchy's square-root ratio cannot
be silently reused for a uniform absolute-value congruence. Rational
or data-adaptive mechanisms are outside this scoped observation.

Absolute-value preconditioning and inertia are established
\cite{vk,vy}; block Krylov, Gaussian range finding, embeddings, and fast
exact rank/characteristic-polynomial algebra are established
\cite{mm,cnw,jps,np}. The new obligations proved here are the finite
relaxed invariant, its noncommuting propagation, and the exact carried
dual lift, together with the unfolded total cost. A global priority
claim requires a separate review.

Principal-component projection through ridge calls begins with
Frostig--Musco--Musco--Sidford \cite{pcp16}. Allen-Zhu--Li's
Theorem~7.3 \cite{pcp17} supplies a gray-band guarantee without an
actual eigengap; Jin--Sidford \cite{pcp19} use rational approximations
and asymmetric systems. These are direct antecedents for the filtering
step and cannot be dismissed by imposing a spectral-gap distinction.
Their ridge or squared-system costs must be charged in the present
separate-factor access model before comparing total runtimes. No
matching full bound from those sources is asserted or ruled out here.
The reviewed gap-free polynomial PCP theorem uses
\(\widetilde O(\gamma^{-1})\) ridge calls. With \(\mu\) denoting
spectral regularization width, the reviewed squared systems have a
\(\mu^2\) positive shift and retain a \(1/\mu\) solve factor
\cite[Theorem 3]{pcp19}. One cannot import \(1/\sqrt\mu\) by treating
that width itself as the strong-convexity modulus. The inspected
reductions do not establish this combined bound. A same-system stochastic
\emph{block} solve can itself batch efficiently; removing a rank factor
from its kernel cost does not supply a global block spectral wrapper.
Stronger compositions remain possible, and the literature audit is not
an exhaustive priority certificate.

The original stronger
\(m=N^{8/9},\theta=1/m\) target remains open: isotropic sketches reach
ambient rows at \(h=m^{-1/16}\), and this hierarchy's displayed full-input
schedule already has exponent at least \(m^{83/32}\), above output.
No statement here bounds all possible factor-access algorithms.

Exact rank tests, Cholesky algebra, interpolation coefficients, and
polynomial recurrences can require arbitrarily large representations.
The carried operators may be nonnormal and their recurrences can grow
exponentially in degree. The proved costs are exact real operation
counts; stable floating-point and bit-complexity guarantees remain separate.

\appendix
\section{An exact-field Gram comparator}\label{app:gram}
The exact-Gram route does not need an eigenvalue or SVD oracle.
Form the two true Grams and use rank revelation and Cholesky to obtain
\(M=Z\mathcal B\), \(Z^\top Z=I\), and explicit positive definite
\(H_c=\mathcal B\mathcal B^\top\) on its product support.
This core algebra costs \(\widetilde O(m^{\omega+\zeta})\), beyond
the two Gram products. If \(r\) reaches the support dimension, return
the exact product as a support projection.

Otherwise use the plain coefficient-ratio bracket and exact shifted
inertia to bisect the positive \(\lambda_{r+1}(H_c)\). In
\(O(\log(m/\theta))\) calls, each costing
\(\widetilde O(m^{\omega+\zeta})\), find an upper endpoint \(u\)
with \(u\leq(1+\theta/100)\lambda_{r+1}(H_c)\).
Draw a fresh Gaussian scalar \(g\) and set
\[
 t=u\left(1+\frac{\theta g^2}{100(1+g^2)}\right).
\]
Then \(K_c=tI-H_c\) is nonsingular almost surely and has at most
\(r\) negative eigenvalues. Draw an independent square Gaussian
\(S\). Every leading principal minor of \(S^\top K_cS\) is nonzero
almost surely: each is a nonzero polynomial in \(S\), as witnessed
by an eigenbasis of the fixed nonsingular symmetric \(K_c\).
Consequently unpivoted recursive symmetric elimination gives
\(S^\top K_cS=LDL^\top\), with diagonal nonzero \(D\), in
\(O(m^{\omega+\zeta})\) field operations. Its Schur-complement
recurrence is \(2T(m/2)+O(m^{\omega+\zeta})\); the leading-minor
condition propagates to each recursive block. Failed zero checks return
within this cap and have probability zero.

Let \(E_-\) select negative diagonal entries and take the physical
space \(\mathcal V=\range(ZS^{-\top}LE_-)\), of width at most \(r\).
For \(x\perp\mathcal V\), \(y=Z^\top x\) satisfies
\(E_-^\top L^\top S^{-1}y=0\). Thus \(y^\top K_cy\geq0\), and
\[
 x^\top MM^\top x=y^\top H_cy\leq t\norm y^2\leq t\norm x^2.
\]
The projection residual is at most
\(\sqrt t\leq(1+\theta/100)\sigma_{r+1}(M)\).
Fast physical range algebra and expansion cost
\(\widetilde O(\Ac(r)+m^{\omega+\zeta})+np+\MM(n,r,p)\).
This proves the quoted Gram comparator in precisely the same arithmetic
model, with no input singular-value-ratio logarithm. It does not assert
numerical stability of exact coefficient signs or Gaussian congruence.

\begin{thebibliography}{99}
\bibitem{companion} \emph{Block spectral approximation of a factored product
from regularized absolute-value pencils}. Companion working manuscript,
November 14, 2026. Product support; finite weighted defect; coefficient
ratio and shifted inertia; safe threshold oracle; filtered Gaussian projection.
\href{absolute-pencil-spectral.pdf}{Local review copy}.
\bibitem{cnw} M. Cohen, J. Nelson, D. Woodruff.
\emph{Optimal Approximate Matrix Product in Terms of Stable Rank}. 2016,
Appendix A.2. \url{https://arxiv.org/abs/1507.02268}.
\bibitem{amanov} A. Amanov, M. van den Berg, J. Zuiddam.
\emph{Multiplying matrices via polynomials: exposition and barriers}. 2026.
\url{https://arxiv.org/abs/2610.30054}.
\bibitem{mm} C. Musco, C. Musco.
\emph{Randomized Block Krylov Methods for Stronger and Faster Approximate
Singular Value Decomposition}. NeurIPS 2015, Algorithm 2 and Theorem 12.
\url{https://arxiv.org/abs/1504.05477}.
\bibitem{np} V. Neiger, C. Pernet.
\emph{Deterministic computation of the characteristic polynomial in the
time of matrix multiplication}. J. Complexity, 2021, Theorem 1.1.
\url{https://arxiv.org/abs/2010.04662}.
\bibitem{jps} C.-P. Jeannerod, C. Pernet, A. Storjohann.
\emph{Rank-profile revealing Gaussian elimination and the CUP matrix
decomposition}. J. Symbolic Computation, 2013, Algorithm 6, Propositions 3, 6.
\url{https://arxiv.org/abs/1112.5717}.
\bibitem{vk} E. Vecharynski, A. V. Knyazev.
\emph{Absolute value preconditioning for symmetric indefinite linear systems}.
SIAM J. Sci. Comput. 35(2), A696--A718, 2013.
\url{https://arxiv.org/abs/1104.4530}.
\bibitem{vy} E. Vecharynski, C. Yang.
\emph{Preconditioned iterative methods for eigenvalue counts}. 2016/2017.
\url{https://arxiv.org/abs/1602.02306}.
\bibitem{pcp16} R. Frostig, C. Musco, C. Musco, A. Sidford.
\emph{Principal Component Projection Without Principal Component Analysis}.
ICML 2016. \url{https://proceedings.mlr.press/v48/frostig16.html}.
\bibitem{pcp17} Z. Allen-Zhu, Y. Li.
\emph{Faster Principal Component Regression and Stable Matrix Chebyshev
Approximation}. ICML 2017, Theorem 7.3.
\url{https://proceedings.mlr.press/v70/allen-zhu17c.html}.
\bibitem{pcp19} Y. Jin, A. Sidford.
\emph{Principal Component Projection and Regression in Nearly Linear Time
through Asymmetric SVRG}. NeurIPS 2019, Definition 1, Theorems 1, 3, and 5.
\url{https://arxiv.org/abs/1910.06517}.
\end{thebibliography}
\end{document}
